\begin{figure}[htbp]
\centering
\begin{tikzpicture}[scale=1.05]
    % Canvas background
    \draw[thick, fill=labbookbg!80!white, draw=labbookrule] (-6.8, -3.0) rectangle (6.8, 3.2);
    \node[anchor=north west, font=\normalsize\bfseries\sffamily, color=labbookslate] at (-6.5, 2.95) {Total Quantum State Space $\mathcal{S}(\ALambda)$ (Convex)};

    % Physical Subspace (Gutzwiller Projection hyperplane)
    \draw[very thick, labbookteal, dashed] (-5.8, 1.6) -- (3.2, 1.6);
    \node[anchor=west, font=\scriptsize\bfseries, color=labbookteal, align=left] at (3.35, 1.6) {Physical Subspace\\$\mathcal{F}_{\mathrm{phys}} = \Ran P_G$ ($\mathcal{D}_i = +\Id$)};

    % Non-Convex Quasi-Free Submanifold M_QF (Shaded Region)
    \draw[very thick, fill=labbookwarm!30!white, draw=labbookcrimson]
        (-4.6, -2.2) [rounded corners=20pt] -- (-2.4, -0.3) [rounded corners=18pt] -- (0.0, -0.9) [rounded corners=18pt] -- (2.0, -0.3) [rounded corners=20pt] -- (4.2, -2.2) [rounded corners=20pt] -- (3.6, -2.4) [rounded corners=20pt] -- (-4.0, -2.4) -- cycle;
    \node[anchor=center, font=\small\bfseries, color=labbookcrimson] at (-0.5, -1.8) {Quasi-Free States $\mathcal{M}_{\mathrm{QF}}$ (Non-Convex)};

    % Pure States: Orthogonal Grassmannian Gr(n, 2n) Boundary
    \draw[very thick, labbookaccent]
        (-4.6, -2.2) to[out=45, in=185] (-2.4, -0.3) to[out=5, in=175] (0.0, -0.9) to[out=-5, in=185] (2.0, -0.3) to[out=5, in=135] (4.2, -2.2);
    \node[anchor=south west, font=\footnotesize\bfseries, color=labbookaccent, align=left] at (-6.2, -0.7) {Grassmannian $\Gr(n, 2n)$\\($\Omega^2 = -\Id$)};

    % State points on Physical Hyperplane
    \filldraw[labbookteal] (-2.5, 1.6) circle (3.5pt) node[anchor=south, font=\small\bfseries, color=labbookteal, yshift=3pt] {$\rho_{\mathrm{exact}}$};
    \filldraw[labbookcrimson] (2.0, 1.6) circle (3.5pt) node[anchor=south, font=\small\bfseries, color=labbookcrimson, yshift=3pt] {$P_G \Omega_{\mathrm{MF}} P_G$};

    % State point on Grassmannian
    \filldraw[labbookaccent] (2.0, -0.3) circle (3.5pt) node[anchor=west, font=\small\bfseries, color=labbookaccent, xshift=6pt] {$\Omega_{\mathrm{MF}} \in \Gr$};

    % Projection Gap Vector (Vertical Arrow)
    \draw[->, very thick, shorten >=5pt, shorten <=5pt, color=labbookslate] (2.0, -0.3) -- (2.0, 1.6);
    \node[anchor=west, font=\scriptsize\bfseries, color=labbookslate, align=left] at (2.15, 0.75) {Gutzwiller Projection\\Gap (BL-01)};

    % Non-Gaussianity Relative Entropy Distance (Dashed Diagonal Arrow)
    \draw[<->, very thick, dashed, shorten >=6pt, shorten <=6pt, color=labbookcrimson] (-2.5, 1.6) -- (2.0, -0.3)
        node[pos=0.42, below=5pt, sloped, font=\scriptsize\bfseries, color=labbookcrimson] {Variational Gap $D_{\mathrm{NG}}(\rho)$ (BL-07)};
\end{tikzpicture}
\caption{Geometry of the variational quantum state space. The full density matrix space $\mathcal{S}(\ALambda)$ is convex. Quasi-free states $\mathcal{M}_{\mathrm{QF}}$ form a non-convex embedded submanifold whose pure boundary is the real orthogonal Grassmannian $\Gr(n, 2n)$. The mean-field state $\Omega_{\mathrm{MF}}$ minimizes energy on the Grassmannian, while Gutzwiller projection maps onto the physical gauge-singlet hyperplane.}
\label{fig:m1_1b_grassmannian_schematic}
\end{figure}
